\chapter{Powder Diffraction}
\label{Chap:powder}

Powder diffraction\index{powder diffraction} was invented by Peter Debye\index{Debye, P.} and Paul Scherrer\index{Scherrer, P.} , as well as independently by Albert Hull\index{Hull, A.} all in 1917, during World War I. 
Working in Germany, Scherrer and Debye were fortunate to be Swiss and (then) Dutch citizens, respectively, as this exempted them from being drafted. Hull worked for a U.S. Company, General Electric, which was a defense contractor.
Having also worked in U.S. companies, he has my sympathies, but I do note that Hull did go on to have several other important inventions. 
Powder diffraction was used initially for the purpose of descriptive crystal-structure determination, for materials where single-crystal specimens for materials were not available. 
It was also initially deployed because the powder diffraction pattern provided a unique structural "fingerprint" for materials.

The name powder diffraction is really a misnomer, as it is used for much more than measurements on powders. A better name for the technique would be polycrystalline diffraction, as it can be used on any sort of bulk material that has a large number of unoriented small crystal components. However, with more than 100 years of tradition behind us, I think it is too late to lobby for a different name. 

An ideal powder diffraction experiment has a large number (typically millions) of crystals that have random orientation. As we will see, powder diffraction is commonly used on samples where the crystal orientation is not completely random. The phenomenon when crystallite orientation ordering is not random is called preferred orientation\index{preferred orientation} and the description of how crystals are oriented is often referred to as texture\index{texture}. 
Powder diffraction is commonly used for the quantification of texture. 
In fact, it is also used for partly- and non-crystalline materials. 
We tend to use the term crystallite\index{crystallite} for each component crystal in a polycrystalline material, where a crystallite is the domain of atoms that all have the same crystallographic order. For metals, the term grain\index{grain} is used for exactly the same meaning as that of crystallite. 

It might be thought that a random orientation of crystallite would produce a random scattering of x-rays, electrons or neutrons, but this is not what happens because diffraction 

As we have seen, diffraction can only occur when the structure factor, $F_{hkl}$, is non-zero and this only occurs for discrete values of $h$, $k$ and $l$ and at specific Q values, where 

$$F_{hkl}=\sum_j^N f_j \exp⁡[-2πi(hx_j + ky_j + lz_j)]$$

Both groups and many others were successful in this work, but as the practice of single-crystal diffraction analysis improved, and structure refinement allowed quantitative structure determina- tion to become routine, one of the key weaknesses of powder diffraction – the loss of information due to reflection overlap – became a major disadvantage. In the relatively few cases where structural refinement from powder diffraction was attempted, reflections that overlapped had to be discarded. Alternately, refinement software was created to treat the specific cases of groups of overlapped reflections for each structural study [for example, see Rietveld (1966)].